%% psl_paper.tex
%% Verification-repair edition, October 2026
%% The pre-repair manuscript is preserved in paper/historical/.
%%
%% Compile with: pdflatex psl_paper.tex

\documentclass[11pt,a4paper]{article}

\usepackage{booktabs}
\usepackage{tabularx}
\usepackage{array}
\usepackage{amsmath}
\usepackage{amssymb}
\usepackage[colorlinks=true,linkcolor=blue,citecolor=blue,urlcolor=blue]{hyperref}
\usepackage[T1]{fontenc}
\usepackage[utf8]{inputenc}
\usepackage{geometry}
\geometry{margin=1in}
\usepackage{authblk}
\usepackage{parskip}

\begin{document}

\title{Structural Constraints for Embedded Control:\\
       A Verification-Repair Study of the P\={a}ninian Systems Language}

\author[1]{Praveen Kumar}
\affil[1]{Independent Researcher, Delhi, India}
\date{October 2026}

\maketitle

\begin{abstract}
The P\={a}ninian Systems Language (PSL) is a research prototype that explores
whether structural ideas inspired by P\={a}nini's grammatical system can be
expressed as useful legality rules for a compact embedded-control language.
PSL lowers programs to a fixed 32-bit instruction format and uses explicit
concepts for context inheritance (Anuvrtti), structured erasure (Lopa),
privilege scope (Adhik\={a}ra), and address visibility (Siddha/Asiddha).

This paper reports both the design and a verification repair undertaken after
an audit found that earlier project claims exceeded the available evidence.
An earlier compile-soundness development had not been exercised by a real
Lean build and relied on a project axiom connecting compile-time and runtime
written-state tracking. The repaired formal model removes that axiom, makes
live written-state, inheritance context, and Adhik\={a}ra depth explicit,
and proves an axiom-free forward simulation from successful canonical control
validation to successful abstract execution. The principal Lean theorems are
typechecked in CI; their dependency audit contains no PSL project axiom and no
\texttt{sorryAx}.

The repaired repository additionally cross-checks exact binaries emitted by the
Python compiler against the Lean validator. The current deterministic corpus
contains 65 accepted word streams (8 checked-in programs and 57 generated
variants), covering 22 distinct ABI words, plus adversarial streams that Lean
must reject. Historical RV32/QEMU runs and a 5,000-program differential test
remain useful experimental evidence, but they are not presented as a proved
refinement from the repaired abstract semantics to hardware.

The result is narrower than the project's earlier claim of a
CompCert-equivalent verified compiler, but it is stronger in trustworthiness:
the formal boundary is explicit, reproducible, and designed to expose rather
than hide remaining gaps.
\end{abstract}

\noindent\textbf{Keywords:}
embedded systems, domain-specific languages, formal methods, Lean 4,
RISC-V, P\={a}nini, privilege separation, state machines

\section{Introduction}

Embedded control software often relies on relationships between operations:
a target must exist before it can be inherited, a value should be written
before it is erased, a privileged operation should occur only inside an
authorized scope, and some operations are intended to remain adjacent or
atomic. Conventional systems languages can express these requirements, but
they frequently encode them through APIs, assertions, static analyzers, code
review, or runtime mechanisms rather than in a small domain-specific syntax.

PSL asks a narrower research question:

\begin{quote}
Can selected control invariants be represented structurally enough that an
invalid instruction sequence is rejected before execution, while keeping the
resulting machine model small enough to audit formally?
\end{quote}

The project borrows terminology and structural inspiration from
P\={a}nini's \emph{A\d{s}\d{t}\={a}dhy\={a}y\={i}}, whose rule system has
long attracted interest from formal-language researchers
\cite{ingerman1967panini,kadvany2016panini}. The historical connection does
not by itself establish novelty or correctness. PSL must be judged by its
actual compiler rules, formal semantics, executable evidence, and comparison
with conventional alternatives.

An October 2026 audit materially changed the project. Earlier documentation
described fifteen proof sprints, a CompCert-style verified compilation chain,
and a Lean theorem analogous to CompCert's
\texttt{compile\_correct} \cite{leroy2009compcert}. A real CI-backed Lean
build showed that this description was not justified: the historical proof
development contained type errors and also depended on a project axiom named
\texttt{p2\_runtime\_correctness}. Several historical examples and firmware
versions also disagreed about core semantics.

Rather than preserve the stronger claim, the repair work narrowed the theorem
and reconciled the language contract. This paper describes that repaired
state.

\subsection{Contributions of the repaired artifact}

The current repository contributes:

\begin{enumerate}
\item A compact 32-bit ABI and Python compiler for a small embedded-control
      DSL with explicit context, erasure, privilege, and scope concepts.
\item A canonical control state that tracks live written targets, live
      Anuvrtti context, and Adhik\={a}ra scope depth.
\item A Lean 4 abstract semantics and an axiom-free theorem that successful
      control validation implies successful abstract execution with the same
      final control state.
\item CI-backed theorem dependency auditing that rejects
      \texttt{sorryAx} and reintroduction of the retired PSL project axiom.
\item A cross-language conformance gate in which exact Python-emitted 32-bit
      words are decoded and checked independently in Lean.
\item An explicit record of semantic contradictions discovered during the
      audit and the decisions used to repair them.
\end{enumerate}

\subsection{Non-claims}

The current artifact does \emph{not} establish:

\begin{itemize}
\item a universal proof that the Python source compiler implements the Lean
      semantics for all possible PSL programs;
\item complete formal semantics for bounded repetition (Avrtti),
      conditional precedence (Utsarga/Apavada), or Sandhi atomicity;
\item a refinement theorem from the Lean machine to a canonical RV32
      firmware implementation;
\item equivalence between QEMU execution and physical RISC-V silicon;
\item that PSL is generally faster, smaller, safer, or superior to existing
      type systems, capability systems, or embedded DSLs.
\end{itemize}

These are research questions, not results.

\section{Paninian-inspired structural concepts}

PSL uses P\={a}ninian terminology as an organizing vocabulary for explicit
compiler state.

\begin{table}[h]
\centering
\caption{PSL concepts and their repaired computational interpretation.}
\begin{tabularx}{\linewidth}{>{\bfseries}p{3.0cm} X}
\toprule
Concept & PSL interpretation \\
\midrule
Karaka & Data-flow role associated with a target or operation. \\
Anuvrtti & Inherit a currently live target context. \\
Avrtti & Encode bounded repetition in the COMP field; full repaired semantics
         remain open. \\
Lopa & Erase live written state and terminate inheritance context. \\
Vrddhi / Ring 0 & Privileged instruction attribute. \\
Siddha / Asiddha & Static address partition at \texttt{0x50}. \\
Adhikara & Explicit lexical/runtime privilege scope. \\
Paribhasha & Named compiler legality rules. \\
Sandhi & Compiler-declared compatible instruction pair; formal atomicity
         remains open. \\
\bottomrule
\end{tabularx}
\end{table}

These mappings should not be interpreted as evidence that the ancient grammar
already contained modern operating-system or compiler mechanisms. They are a
design analogy whose usefulness must be demonstrated by the resulting formal
and executable properties.

\section{Instruction representation}

Every lowered PSL instruction is one 32-bit big-endian word:

\[
\begin{array}{c|c|c|c|c|c}
31\ldots28 & 27\ldots24 & 23\ldots16 & 15\ldots8 & 7\ldots4 & 3\ldots0 \\
\hline
\mathrm{RING} & \mathrm{COMP} & \mathrm{OPCODE} &
\mathrm{TARGET} & \mathrm{FLAGS} & \mathrm{COND}
\end{array}
\]

The repaired core recognizes:

\begin{center}
\begin{tabular}{lll}
\toprule
Value & Name & Current control meaning \\
\midrule
\texttt{0x00} & Lopa & structured erasure \\
\texttt{0x05} & Write & write-class operation \\
\texttt{0x06} & Read & read operation \\
\texttt{0xCC} & Store & privileged write-class operation \\
\texttt{0xAA} & Adhikara OPEN & increment scope depth \\
\texttt{0xBB} & Adhikara CLOSE & decrement non-zero scope depth \\
\bottomrule
\end{tabular}
\end{center}

\texttt{COMP=1} denotes Anuvrtti. Values greater than one are historically
used for Avrtti repetition counts but are not yet given full execution
semantics by the repaired Lean theorem.

Ring is an instruction attribute. The audit removed the earlier formal design
in which the machine also stored a mutable global ring changed by OPEN/CLOSE,
because that duplicated privilege authority and produced a false theorem about
restoring arbitrary prior ring state.

\section{Canonical control semantics}

\subsection{State}

The repaired control state is

\[
C = \langle W, K, D \rangle
\]

where:

\begin{itemize}
\item $W : \mathrm{Addr} \rightarrow \mathrm{Bool}$ records whether an
      address currently contains live written state;
\item $K : \mathrm{Option}\;\mathrm{Addr}$ is the live Anuvrtti target;
\item $D : \mathbb{N}$ is the active Adhik\={a}ra scope depth.
\end{itemize}

Initially, $W(a)=\mathrm{false}$ for all addresses, $K=\mathrm{none}$, and
$D=0$.

For an ABI word $w$,

\[
\mathrm{resolve}(C,w) =
\begin{cases}
K & \text{if } w.\mathrm{COMP}=1 \\
\mathrm{some}(w.\mathrm{TARGET}) & \text{otherwise.}
\end{cases}
\]

An Anuvrtti instruction is invalid if $K=\mathrm{none}$.

\subsection{Address and privilege rules}

The address partition is:

\[
\mathrm{Siddha}(a) \Leftrightarrow a < \texttt{0x50},
\qquad
\mathrm{Asiddha}(a) \Leftrightarrow a \ge \texttt{0x50}.
\]

The repaired rules are:

\begin{enumerate}
\item Any Ring-0 executable instruction must resolve to Asiddha.
\item Any Ring-0 executable instruction requires $D>0$.
\item Store always requires Ring 0 and $D>0$.
\item Lopa of an Asiddha target requires Ring 0 and $D>0$.
\item Ring-2 Write to Asiddha is legal. Its visibility consequences belong to
      the future memory/refinement model rather than the control validator.
\end{enumerate}

The last item repairs a direct inconsistency in the historical artifact:
Sprint-9 compiler/firmware examples allowed a Ring-2 Write to Asiddha, while
the old Lean machine rejected it.

\subsection{Live written-state and Lopa}

Write and Store establish live state:

\[
W[t] := \mathrm{true}.
\]

Lopa requires $W[t]=\mathrm{true}$ and then consumes it:

\[
W[t] := \mathrm{false}.
\]

Therefore a second Lopa on the same target requires a new intervening
Write/Store.

Lopa also clears $K$. This resolves a second historical contradiction.
Sprint 5 defined Lopa as an inheritance boundary, while the later Sprint-13
key-lifecycle example attempted an inherited Lopa immediately after an
explicit Lopa. The repaired language adopts the earlier, structurally cleaner
meaning: inheritance cannot cross erasure.

\subsection{Adhikara}

OPEN and CLOSE update scope depth:

\[
\mathrm{OPEN}: D := D+1,
\qquad
\mathrm{CLOSE}: D>0;\; D := D-1.
\]

A complete valid program must terminate with $D=0$. A CLOSE at depth zero is
invalid.

\section{Compiler architecture}

The Python compiler performs symbol resolution, IR construction, Sandhi
annotation, structural validation, lowering checks, and final ABI emission.
The repaired validator independently tracks live written targets, live context,
and scope depth instead of assuming that IR construction already enforced the
rules correctly.

The active examples are regression-tested against exact expected ABI words.
Four examples that crossed semantic eras were repaired:

\begin{itemize}
\item \texttt{lopa\_core.pvm} now terminates at the Lopa boundary;
\item \texttt{sandhi\_core.pvm} uses a scoped same-ring Store/Store pair;
\item \texttt{paribhasha\_valid\_core.pvm} scopes its privileged Store;
\item \texttt{key\_lifecycle.pvm} removes the inherited Lopa through an
      erasure boundary.
\end{itemize}

Their pre-repair text is preserved under \texttt{programs/historical/}.
Historical sprint evidence is provenance, not automatically the current
language definition.

\section{Lean model and proved result}

The machine-checked authority is
\texttt{lean/PSL/Semantics.lean}, typechecked under pinned Lean 4.34.1
\cite{moura2021lean4}.

The abstract machine contains a control state and a simple memory function.
The memory model intentionally abstracts away a full hardware-level
Siddha/Asiddha implementation. A step first executes the canonical control
transition. Only a successful control transition permits the abstract memory
update.

Two folds are defined over a word stream:

\begin{itemize}
\item \texttt{validateWords}: applies only the control transition;
\item \texttt{execute}: applies control plus abstract memory transition.
\end{itemize}

The central theorem is \texttt{validateWords\_execution}. Informally:

\[
\begin{split}
& s_0.\mathrm{control}=c_0 \;\land\;
\mathrm{validateWords}(ws,c_0)=\mathrm{some}(c_f) \\
& \qquad \Rightarrow
\exists s_f.\;
\mathrm{execute}(ws,s_0)=\mathrm{some}(s_f)
\;\land\;
s_f.\mathrm{control}=c_f.
\end{split}
\]

Corollaries connect successful validation of lowered IR to abstract execution
and require closed Adhik\={a}ra scope for the definition of a
\texttt{ValidProgram}.

This theorem is intentionally not called a complete compiler-correctness
theorem. Validation and execution share a canonical control transition, so
the result establishes consistency between the validator and the abstract
machine. It does not by itself formalize the Python parser or prove every
source transformation correct.

The theorem dependency audit currently reports only Lean's standard
\texttt{propext}; there is no PSL project axiom and no \texttt{sorryAx}.
The historical \texttt{p2\_runtime\_correctness} axiom has been removed.

\section{Cross-language conformance evidence}

A practical boundary remains between the Python compiler and the Lean model.
The current artifact addresses it with finite, reproducible cross-language
conformance rather than pretending it is a universal proof.

CI performs the following sequence:

\begin{enumerate}
\item compile eight checked-in canonical PSL programs using Python;
\item generate additional deterministic valid sources from the currently
      formalized subset;
\item collect exact numeric 32-bit ABI words;
\item generate Lean propositions requiring every word stream to satisfy
      \texttt{validateEncodedClosed};
\item require every distinct emitted word to round-trip through Lean's
      numeric ABI decoder;
\item typecheck these propositions with Lean;
\item typecheck a separate adversarial file whose malformed streams must be
      rejected.
\end{enumerate}

At the time of this repair, the passing positive corpus contains 65 word
streams: eight checked-in programs plus 57 generated variants, covering 22
distinct ABI words.

Negative cases include:

\begin{itemize}
\item Anuvrtti without live context;
\item CLOSE at zero scope depth;
\item Ring-0 access to Siddha;
\item Store outside Adhik\={a}ra;
\item Ring-2 Store even inside Adhik\={a}ra;
\item Asiddha Lopa without a live prior write;
\item double Lopa without a new write;
\item post-Lopa Anuvrtti;
\item unclosed Adhik\={a}ra scope.
\end{itemize}

This evidence is stronger than the project's historical theorem-name string
checks, but it remains finite.

\section{Historical experimental evidence}

Earlier project sprints produced freestanding RV32 firmware executed under
QEMU and a 5,000-program differential test run. These artifacts remain useful
for regression archaeology and for understanding how the language evolved.

They are now classified as \emph{historical executable evidence} for three
reasons. First, firmware semantics changed between sprints; there is not yet
one canonical RV32 runtime. Second, some historical examples contradict the
repaired language contract. Third, QEMU execution is not a proof about
physical silicon.

The 5,000-program differential experiment is likewise useful but should not
be read as formal verification. It compared implementations/properties over a
generated domain and helped discover an overconstrained Sandhi rule. Its
scientific value is empirical falsification, not theorem proving.

\section{Sandhi boundary}

Sandhi remains an interesting but deliberately open part of the repaired
formal story. The compiler checks pair compatibility and marks the first word
with \texttt{FLAGS=0xE}. Historical firmware interpreted this marker as a
request for paired execution.

The current Lean theorem does not prove atomicity, interruption behavior,
rollback behavior, or hardware visibility of such a pair. Therefore this
paper uses the phrase \emph{Sandhi pair marking}, not proved atomic execution.
A proper result requires an explicit atomic transition model and a refinement
to the selected RV32 runtime.

\section{Threats to validity and limitations}

\paragraph{Shared transition definition.}
The validator-to-execution theorem is a consistency result over a shared
canonical control transition. It should not be confused with an independently
verified optimizing compiler.

\paragraph{Finite compiler bridge.}
The 65-stream cross-language corpus can find disagreements but cannot prove
the absence of disagreements outside the generated domain.

\paragraph{Incomplete language features.}
Avrtti, Utsarga/Apavada, and Sandhi atomicity are not yet incorporated into
the central repaired theorem.

\paragraph{Simplified memory model.}
The Lean memory function does not yet model realistic MMIO ordering, caches,
interrupts, DMA, speculation, or a concrete Siddha/Asiddha hardware
implementation.

\paragraph{Historical firmware plurality.}
Different RV32 sprint firmware files implement different subsets or semantic
eras. A single canonical runtime must be selected before an implementation
refinement claim is meaningful.

\paragraph{No physical-hardware proof.}
QEMU evidence does not establish physical RISC-V behavior.

\paragraph{Comparative novelty.}
The use of P\={a}ninian terminology does not by itself establish a novel
formal-methods result. Any novelty claim requires comparison with existing
effect systems, typestate, capability machines, verified DSLs, protocol state
machines, and certified compilation.

\section{Future work}

The repair establishes an order for future work:

\begin{enumerate}
\item formalize a source/IR relation or proof-producing certificate that
      reduces the remaining Python-to-Lean gap;
\item add independent formal semantics for Avrtti and conditional execution;
\item give Sandhi an explicit atomic transition semantics and adversarial
      interruption model;
\item select one canonical RV32 runtime;
\item establish transition-level differential/refinement evidence between
      that runtime and Lean;
\item run on physical hardware;
\item only after those steps, evaluate whether a stronger end-to-end
      compiler-correctness theorem is justified.
\end{enumerate}

Performance, code-size, energy, and determinism claims should be evaluated
against explicit conventional baselines rather than inferred from language
structure.

\section{Conclusion}

PSL remains an experimental systems-language project, but the verification
repair changes what can responsibly be said about it.

The strongest current formal statement is not that PSL is a fully verified
compiler or that P\={a}ninian structure has been proved superior to modern
systems techniques. The result is narrower: the repaired canonical control
validator is machine-checked in Lean, successful validation is proved to
forward-simulate into successful abstract execution, the theorem uses no PSL
project axiom, and the current Python compiler is continuously cross-checked
against that model over a deterministic positive and negative corpus.

The audit also produced a methodological result: counts of sprints, checks,
and green scripts are poor substitutes for an explicit trust boundary. Once
the repository began requiring a real Lean build and distinguishing tests from
proofs, several contradictions became visible and repairable.

That narrower result provides a credible base for the next question: whether
the same structural contract can survive refinement into one real RV32
runtime and, eventually, physical hardware.

\section*{Acknowledgements}

The author thanks the Lean, QEMU, and GNU/RISC-V communities whose open-source
infrastructure makes independent reproduction possible.

\bibliographystyle{plain}
\bibliography{psl_paper}

\end{document}
